% Generated from locale/tr/content/set-theory/card-arithmetic/opps.tex; do not hand-edit.


\olfileid[tr]{sth}{card-arithmetic}{opps}
\olsection{Temel İşlemlerin Tanımlanması}

Sırayı izlememiz gerekmediğinden kardinal aritmetiğini tanımlamak sıra
sayısı aritmetiğini tanımlamaktan oldukça kolaydır. Toplama, çarpma ve
üs almayı aynı anda tanımlayacağız.

\begin{defn}
$\cardfont{a}$ ve $\cardfont{b}$ kardinal sayılar olduğunda:
\begin{align*}
	\cardfont{a} \cardplus \cardfont{b} &\defis
	\card{\cardfont{a} \disjointsum \cardfont{b}}\\
	\cardfont{a} \cardtimes \cardfont{b} &\defis
	\card{\cardfont{a} \times \cardfont{b}}\\
	\cardexpo{\cardfont{a}}{\cardfont{b}} &\defis
	\card{\funfromto{\cardfont{b}}{\cardfont{a}}}
\end{align*}
burada $\funfromto{X}{Y} = \Setabs{f}{f \text{ bir } X \to Y
\text{ fonksiyonudur}}$ olur. (Her $X$ ve $Y$ kümesi için
$\funfromto{X}{Y}$'nin var olduğunu göstermek kolaydır; bunu alıştırma
olarak bırakıyoruz.)
\end{defn}

\begin{prob}
Her $X$ ve $Y$ kümesi için $\funfromto{X}{Y}$'nin var olduğunu $\Zminus$
içinde ispatlayın. $\ZF$ içinde çalışarak
$\setrank{\funfromto{X}{Y}}$ değerini, $\setrank{X}$ ve $\setrank{Y}$
değerlerinden
\olref[sth][ord-arithmetic][using-addition]{rankcomputation} tarzında
hesaplayın.
\end{prob}

Bu tanımı açıklamak yararlı olabilir. Toplama bakımından tanım,
\olref[ord-arithmetic][add]{defdissum} içinde tanımlanan ayrık toplam
$\disjointsum$ kavramını kullanır; sonlu durumlarda doğru sonucu verdiğini
görmek kolaydır. Çarpma bakımından
\olref[sfr][set][pai]{cardnmprod}, $A$'nın $n$ elemanı ve $B$'nin $m$
elemanı varsa $A\times B$'nin $n\cdot m$ elemanı olduğunu söyler;
dolayısıyla tanımımız düşünceyi yalnızca sonlu ötesi çarpmaya geneller.
Üs alma da benzerdir: sonlu durumdaki düşünceyi sonlu ötesi duruma
genelliyoruz. Gerçekten sonlu ötesi kardinal aritmetiği bazı yönlerden
sonlu ötesi sıra sayısı aritmetiğine kıyasla ``olağan'' aritmetiğe çok
daha fazla benzer:

\begin{prop}\ollabel{cardplustimescommute}
$\cardplus$ ve $\cardtimes$ değişmeli ve birleşmelidir.
\end{prop}

\begin{proof}
Değişmelilik için
\olref[cardinals][cardsasords]{lem:CardinalsBehaveRight} gereği
$\cardeq{(\cardfont{a} \disjointsum
\cardfont{b})}{(\cardfont{b} \disjointsum \cardfont{a})}$ ve
$\cardeq{(\cardfont{a} \times \cardfont{b})}{(\cardfont{b} \times
\cardfont{a})}$ olduğunu gözlemlemek yeterlidir. Birleşmeliliği alıştırma
olarak bırakıyoruz.
\end{proof}

\begin{prob}
$\cardplus$ ve $\cardtimes$ işlemlerinin birleşmeli olduğunu ispatlayın.
\end{prob}

\begin{prop}
$A$ sonsuzdur ancak ve ancak
$\card{A} \cardplus 1 = 1 \cardplus \card{A} = \card{A}$.
\end{prop}

\begin{proof}
\olref[cardinals][classing]{generalinfinitycharacter} içindeki gibi;
\olref[ord-arithmetic][using-addition]{ordinfinitycharacter} ve
\olref[cardinals][cardsasords]{lem:CardinalsBehaveRight} sonuçlarından
çıkar.
\end{proof}

Bu, sıra sayısı toplaması/çarpması ile kardinal toplaması/çarpması için
niçin farklı simgeler kullanmamız gerektiğini açıklar: bunlar gerçekten
\emph{farklı} işlemlerdir. Sıradaki iki sonuç, sıra sayısı üs alması ile
kardinal üs almasının da farklı işlemler olduğunu gösterir.
(\olref[z][infinity-again]{defnomega} sonucunun $2=\{0,1\}$ gerektirdiğini
hatırlayın):

\begin{lem}\ollabel{lem:SizePowerset2Exp}
Her $A$ için $\card{\Pow{A}} = \cardexpo{2}{\card{A}}$.
\end{lem}

\begin{proof}
Her $B\subseteq A$ altkümesi için $\chi_B\in\funfromto{A}{2}$ fonksiyonu
şöyle verilsin:
\begin{align*}
	\chi_{B}(x) &\defis
	\begin{cases}
		1 & \text{eğer }x\in B\\
		0 & \text{aksi hâlde.}
	\end{cases}
\end{align*}
Şimdi $f(B)=\chi_B$ olsun; bu, $f\colon\Pow{A}\to
\funfromto{A}{2}$ !!a{bijection} tanımlar. Öyleyse
$\cardeq{\Pow{A}}{\funfromto{A}{2}}$. Dolayısıyla
$\cardeq{\Pow{A}}{\funfromto{\card{A}}{2}}$ ve böylece
$\card{\Pow{A}} = \card{\funfromto{\card{A}}{2}} =
2^{\card{A}}$ olur.
\end{proof}

Bu kısa ispat esasen \olref[sfr][siz][red-alt]{sec} tartışmasını kapsar.
Orada $\Pow{\omega}$'nın sayılamazlığının sonsuz ikili diziler kümesi
$\Bin^\omega$'nın sayılamazlığına nasıl ``indirgeneceğini'' gösterdik.
Aslında $\Bin^\omega$ yalnızca $\funfromto{\omega}{2}$'dir; önceki ispat
da \olref[sfr][siz][red-alt]{sec} içindeki uslamlamanın $\omega$ yerine
herhangi bir $A$ kümesi konarak yürütülebileceğini gösterdi. Sonuç ayrıca
kardinal üs alma hakkında hızlı bir olgu verir:

\begin{cor}\ollabel{cantorcor}
Her $\cardfont{a}$ kardinal sayısı için
$\cardfont{a}<\cardexpo{2}{\cardfont{a}}$.
\end{cor}

\begin{proof}
Cantor Teoremi'nden (\olref[sfr][siz][car]{thm:cantor}) ve
\olref{lem:SizePowerset2Exp} sonucundan çıkar.
\end{proof}
\noindent
Dolayısıyla $\omega<\cardexpo{2}{\omega}$. Fakat dikkat edin: bu,
\emph{kardinal} üs alma hakkındaki bir sonuçtur. \emph{Sıra sayısı} üs
almasıyla karşılaştırılmalıdır; çünkü ikinci durumda
$\omega=\ordexpo{2}{\omega}$ olur (bkz.
\olref[ord-arithmetic][expo]{sec}).

Kardinal üs alma konusundayken $\Real$'in hangi ``anlamda''
!!{nonenumerable} olduğu konusunda biraz daha kesin konuşabiliriz.

\begin{thm}\ollabel{continuumis2aleph0}
$\card{\Real} = \cardexpo{2}{\omega}$
\end{thm}

\begin{proof}[İspat taslağı]
Bunu ispatlamanın birçok yolu vardır. En dolaysız yol,
$\cardle{\Pow{\omega}}{\Real}$ ve $\cardle{\Real}{\Pow{\omega}}$
olduğunu savunmak; sonra Schröder--Bernstein ile
$\cardeq{\Real}{\Pow{\omega}}$ ve
\olref[card-arithmetic][opps]{lem:SizePowerset2Exp} ile
$\card{\Real}=\cardexpo{2}{\omega}$ sonucunu çıkarmaktır.
$f\colon\Pow{\omega}\to\Real$ ve $g\colon\Real\to\Pow{\omega}$
iğnelemeleri tanımlamayı (aydınlatıcı) bir alıştırma olarak bırakıyoruz.
\end{proof}

\begin{prob}
$\cardle{\Pow{\omega}}{\Real}$ ve $\cardle{\Real}{\Pow{\omega}}$
olduğunu göstererek
\olref[sth][card-arithmetic][opps]{continuumis2aleph0} ispatını
tamamlayın.
\end{prob}

